% Linux KVM layering: host processes and a QEMU-KVM guest above a
% Linux kernel with the KVM module, on VT-x/AMD-V hardware. Clean
% remake of the talk's slide-18 diagram. Used in M12-L03.
\documentclass[tikz,border=4pt]{standalone}
\usepackage{tikz}
\input{_m10-styles.texinput}
\begin{document}
\begin{tikzpicture}[m10, x=1mm, y=-1mm]
  % Host userspace processes
  \node[greybox, minimum width=22mm, minimum height=30mm,
        font=\sffamily, align=center, anchor=north west]
        at (0,8) {host\\process};
  \node[greybox, minimum width=22mm, minimum height=30mm,
        font=\sffamily, align=center, anchor=north west]
        at (25,8) {host\\process};
  \node[note, font=\Large, text=labelgrey] at (55.5,23) {$\cdots$};
  % The VM: a dashed container holding the guest stack
  \draw[draw=stackedge, dashed, line width=0.8pt, rounded corners=2.5pt]
        (63,0) rectangle (106,40);
  \node[note, anchor=south west, font=\sffamily\small,
        text=stackedge] at (63,-1) {one VM};
  \node[safebox, minimum width=37mm, minimum height=9mm,
        font=\sffamily, anchor=north west] at (66,3)
        {guest userspace};
  \node[safebox, minimum width=37mm, minimum height=9mm,
        font=\sffamily, anchor=north west] at (66,15)
        {guest kernel};
  \node[warnbox, minimum width=37mm, minimum height=9mm,
        font=\sffamily, anchor=north west] at (66,27)
        {QEMU-KVM};
  % Host Linux kernel with KVM module
  \draw[rounded corners=2.5pt, fill=stackfill, draw=stackedge,
        line width=0.7pt] (0,44) rectangle (106,58);
  \node[note, font=\large\bfseries, text=stackedge] at (38,51)
        {Linux kernel};
  \node[box, fill=white, draw=stackedge, font=\sffamily,
        minimum height=8mm, anchor=west] at (72,51) {KVM module};
  % Hardware
  \draw[rounded corners=2.5pt, fill=labelgrey!12, draw=labelgrey!70,
        line width=0.7pt] (0,62) rectangle (106,74);
  \node[note, font=\large\bfseries, text=labelgrey] at (53,68)
        {Hardware (VT-x / AMD-V)};
\end{tikzpicture}
\end{document}
